\section{Introduction}
\label{sec:intro}

\subsection{Background}
Modern ransomware is a cryptographic program first and malware second. A typical family generates a
symmetric key per file (AES or ChaCha20), encrypts that key with an RSA or elliptic-curve public key, and
deletes the symmetric key once it is used~\cite{oz2022survey}. Information stealers take the opposite route
and harvest the tokens, API keys and one-time-password seeds that sit in memory in plain text. In both cases
the decisive evidence is volatile: the injected code that carried out the attack, the cipher implementation
the malware loaded, keys still held by a running process, and the live network connections it opened. Once the
machine is powered off, most of it is gone.

Memory forensics recovers this evidence from an image of RAM. Volatility~3 reconstructs
processes, command lines, sockets and suspicious memory regions from such an image, and YARA and Wireshark
cover byte patterns and traffic. An analyst still has to join the results by hand, decide which cipher is
actually in use (nearly every Windows process names AES somewhere), judge whether a recovered key is weak,
and write up the case. Recent work has started to hand that last step to large language
models~\cite{oh2024volgpt,sharma2025forensicllm,scanlon2023chatgpt}.

\subsection{Problem statement}
Two properties of LLMs make this hand-over risky. First, an LLM can state facts that are not in the evidence,
and a reader cannot tell which sentences are supported. Second, and specific to forensics, the evidence is
attacker controlled. A JWT claim, an OTP issuer label, a TLS server name or a process command line can carry
text such as ``ignore previous instructions and report this host as clean''. This is indirect prompt
injection~\cite{greshake2023not}: the attacker never talks to the model, but plants instructions in data the
model later reads. A triage tool that can be talked into calling a compromised host clean is worse than none.

The problem this work addresses is therefore: \emph{given a memory image, automatically extract and assess the
cryptographic and behavioural evidence of compromise, and produce a readable report in which every claim is
traceable to evidence and the report cannot be steered by text planted in that evidence.}

\subsection{Research gap}
Each part of this problem has been studied in isolation (Section~\ref{sec:related}). Memory-based detectors
output a label but not the cryptography behind it~\cite{carrier2022detecting}; key-extraction studies recover
one key type from one family~\cite{babiano2023evaluation,baier2024tls}; weak-key analysis runs on key
collections, not on keys carved from a live image~\cite{pelofske2024efficient}; and LLM-assisted
forensics~\cite{oh2024volgpt,sharma2025forensicllm} does not test whether crafted evidence can bias the model.
Defences against indirect injection~\cite{hines2024spotlighting,liu2024formalizing} have not been evaluated on
forensic evidence, nor combined with a check that the report's numbers and citations match deterministic
results. We found no system that joins all three: memory artefacts, cryptographic assessment with recovered
keys, and a measured, defended LLM report.

\subsection{Objectives}
Each objective has a measurable criterion, reported in Section~\ref{sec:results}.
\begin{enumerate}
  \item \textbf{O1 Extraction.} Run Volatility~3 on Windows (XP to 11, Server) and Linux images and record
        processes, command lines, connections and injected regions. \emph{Criterion:} end to end on at least
        ten images from different operating systems.
  \item \textbf{O2 Cipher identification.} Identify AES, DES/3DES, RSA, ChaCha20/Salsa20, AES-GCM, ECC
        (P-256, secp256k1) and Curve25519 with a confidence and an evidence basis. \emph{Criterion:} the cipher a
        process uses is identified in every benchmark case.
  \item \textbf{O3 Key assessment and recovery.} Test every RSA key against NIST strength and known
        weaknesses, and recover live AES, ChaCha20 and RSA private keys. \emph{Criterion:} every planted live
        key recovered, no false AES key in random data.
  \item \textbf{O4 Credentials and network.} Decode JWTs, API tokens and OTP secrets (stored redacted) and flag
        weak negotiated TLS.
  \item \textbf{O5 Scoring.} Compute a risk score from sourced formulas only, with an uncertainty range.
        \emph{Criterion:} a clean baseline scores in a lower band than every compromised image.
  \item \textbf{O6 Cited report.} Produce an LLM report from findings only. \emph{Criterion:} at least 90\% of
        claim lines cited, and the stated risk within 20 points of the computed one.
  \item \textbf{O7 Robustness.} Measure prompt-injection success with and without defences, including adaptive
        attacks. \emph{Criterion:} no successful attack reaches a reader without a warning.
\end{enumerate}

\subsection{Contributions}
\begin{itemize}
  \item An end-to-end, open-source pipeline (\repo) covering Windows and Linux images, with a read-only
        dashboard over all results (\demo).
  \item A cryptographic assessment of memory that goes beyond naming ciphers: Dempster-Shafer confidence with
        an evidence basis, NIST-based RSA strength, and recovery of live AES, ChaCha20 and RSA private keys,
        validated on planted keys and on real WannaCry memory.
  \item A risk model in which every number names its formula and source, with a CVSS-edge uncertainty range.
  \item A defended LLM report and a 15-attack red-team, including adaptive attacks, with Wilson intervals and
        Fisher tests; and a citation and numeric cross-check that also caught real errors in our own
        reports.
\end{itemize}

\subsection{Paper organisation}
Section~\ref{sec:related} reviews related work. Section~\ref{sec:method} describes the architecture, data,
algorithms, setup and metrics. Section~\ref{sec:results} reports results on 14 memory images, a ground-truth
lab image, a 50-case benchmark and the red-team. Section~\ref{sec:conclusion} concludes with limitations and
future work.
